\clearpage
\section{Overview of the Monochromator Control Electronics}
\label{Overview of the Monochromator Control Electronics}

This appendix briefly describes the general layout of the electronic control of
the monochromator's servo motors, which set the reflection grating and allow to
tune the monochromator to the desired wavelength. The details of all the
electronics are described in the accompanying manual, however there are some
minor differences between the manual and the actual implementation.

The purpose of the entire system is to allow control of the servo motors
from a computer, and to manage the power supply.
Fig.~\ref{image:Schaltplan_Leistungselektronik} is a schematic of the entire
circuit layout in all its major parts. The two central components are the TSNM
boards, one for each motor, which contain the actual control electronics. Each
of them receives input from the computer on the pins labeled 10, 24, and 23, and
returns signals via pin 12 and 10 to the controlling speedometers (T1, T2),
which control the speed of the motors (M1, M2). Also, from the connectors A1 and
A2, the TSNM gives power to the motors through limiter coils L1/L2, which
prevent power surges. Missing from the schematic is the PSNM board, which
regulates the power for both the TSNM's. There is a separate system, the
\emph{Heidenhain} system, also not indicated, that monitors the servo motors and
provides feedback to the electronics.

Figs.~\ref{image:connection_panel} -- \ref{image:Schaltplan_Leistungselektronik}
contain photographs of the system's components in their physical implementation.

\begin{figure}[htbp]
    \centering
    \includegraphics[width=\textwidth]{images/connection_panel.eps}
    \caption{Connection Panel}
  \label{image:connection_panel}
\end{figure}

Fig.~\ref{image:connection_panel} shows the front panel of the entire system.
From the user's perspective, this panel combines input and output. The
computer (which is running an appropriate software) is plugged in at the
\emph{Bus System} circuit board. The bus has to be connected to the lower panel,
along the bright green lines. The \emph{Main Power Supply} is where the entire
system receives power. The \emph{Servo Power Supply} and \emph{Servo Control
Connectors} constitute the system's output. The power supply has to be connected
directly to the motors, in Fig.~\ref{image:Schaltplan_Leistungselektronik} this
missing connection would be the lines between the limiter coils L1/L2 and the
motors M1/M2. The \emph{Servo Control Connectors} have to be connected to the
controlling speedometers T1/T2, which are physically attached to the motors.
This missing connection is the one labeled \emph{Control Circuit} in
Fig.~\ref{image:Schaltplan_Leistungselektronik}. The Heidenhain control
system is also plugged into this panel.

\begin{figure}[htbp]
    \centering
\includegraphics[width=\textwidth]{images/tsnm_psnm_back.eps}
    \caption{PSNM/TSNM Board}
  \label{image:tsnm_psnm_back}
\end{figure}

\begin{figure}[htbp]
    \centering
\includegraphics[width=0.9\textwidth]{images/tsnm_psnm_front.eps}
    \caption{Frontside of PSNM/TSNM Board}
  \label{image:tsnm_psnm_front}
\end{figure}

Fig.~\ref{image:tsnm_psnm_back} shows the system's central component, the
PSNM/TSNM board. The PSNM receives power from the three blue cables, which
originate from the \emph{Main Power Supply}
(Fig.~\ref{image:transformator_throttle}). It relays power to the two TSNM's and
indirectly to the servo motors (the connection to the servo motors is the gray
cables from the A1/A2 connectors of the TSNM's). The power connection between
PSNM and TSNM is visible as the two dark green bars on the board. The \emph{TSNM
Control Circuit and Control Unit Connector} are in the lower
half of the board. The ribbon cable connects to the \emph{Connection Panel}
(Fig.~\ref{image:connection_panel}). In
Fig.~\ref{image:Schaltplan_Leistungselektronik}, the ribbon cable would connect
between the pins 10-24 of both TSNM's and the T1/T2 speedometers (\emph{Control
Circuit}) and the computer (\emph{Control Unit}), respectively. On the side of
the \emph{Connection Panel}, the ribbon cable splits up and connects to the four
outlets of the \emph{Control Unit} and \emph{Servo Control Connectors}
(Fig.~\ref{image:connection_panel}). The brown cables belong to the Heidenhain
system.

Fig.~\ref{image:tsnm_psnm_front} shows the front side of the PSNM/TSNM board.
The PSNM and the two TSNM's are plug-in cards.

\begin{figure}[htbp]
    \centering
\includegraphics[width=\textwidth]{images/transformator_throttle.eps}
    \caption{Power Transformer and Limiter Coils}
  \label{image:transformator_throttle}
\end{figure}

Fig.~\ref{image:transformator_throttle} shows the intermediate parts between
the \emph{Connection Panel} (Fig.~\ref{image:connection_panel}) and the
\emph{PSNM/TSNM Board} (Fig.~\ref{image:tsnm_psnm_back}). First, this is the
\emph{Main Power Supply} that distributes the power to all components. It draws
power from the \emph{Main Power Supply} plug on the \emph{Connection Panel}
(Fig.~\ref{image:connection_panel}) via the black cable at the top. The three
bundled blue cables go from the \emph{Main Power Supply} to the PSNM
(Fig.~\ref{image:tsnm_psnm_back}). Additionally, there is power going to the
Heidenhain system, through the device labled as \emph{MET-O-PHM}.

Second, there are the four \emph{Limiter Coils}, which in
Fig.~\ref{image:Schaltplan_Leistungselektronik} are labeled L1/L2. The have the
bundled gray cables as input, which originate from the A1/A2 connectors on the
TSNM's (Fig.~\ref{image:tsnm_psnm_back}), and continue with two gray and two
blue cables which go the \emph{Servo Power Supply} on the \emph{Connection
Panel} (Fig.~\ref{image:connection_panel}). The purpose of the coils is simply
to limit power surges that would occur when the motors switch on or off.

Fig.~\ref{image:transformator_throttle} also shows the backside of the
\emph{Connection Panel} (Fig.~\ref{image:connection_panel}) on the very left.
The splitting-up of the ribbon cable into the \emph{Control Unit} and the
\emph{Servo Control Connectors} is visible.

\begin{figure}[htbp]
    \centering
\includegraphics[width=\textwidth]{images/Schaltplan_Leistungselektronik.eps}
    \caption{Circuit Layout of Monochromator Control Electronics}
  \label{image:Schaltplan_Leistungselektronik}
\end{figure}

